\documentclass[11pt]{article}
\usepackage[margin=1.0in]{geometry}
\usepackage{amsmath,amssymb,amsthm,mathtools}
\usepackage[T1]{fontenc}
\usepackage[expansion=false]{microtype}
\usepackage{enumitem}
\usepackage{booktabs}

% ---------- theorem environments ----------
\newtheorem{theorem}{Theorem}
\newtheorem{lemma}{Lemma}
\theoremstyle{definition}
\newtheorem{requirement}{Structural requirement}
\newtheorem*{definition}{Definition}
\newtheorem*{semantics}{Semantic judgment}
\newtheorem*{assumption}{Additional assumption}
\theoremstyle{remark}
\newtheorem*{remark}{Remark}

% ---------- notation ----------
\newcommand{\fin}{\subseteq_{\mathrm{fin}}}
\newcommand{\Past}{\mathsf{Past}}
\newcommand{\Auth}{\mathsf{Auth}}
\newcommand{\Permit}{\mathsf{Permit}}
\newcommand{\Due}{\mathsf{Due}}
\newcommand{\NewDue}{\mathsf{NewDue}}
\newcommand{\Resolve}{\mathsf{Resolve}}
\newcommand{\Continue}{\mathsf{Continue}}
\newcommand{\Met}{\mathsf{Met}}
\newcommand{\AddPre}{\mathsf{AddPre}}
\newcommand{\DropPre}{\mathsf{DropPre}}
\newcommand{\Designate}{\mathsf{Designate}}
\newcommand{\Pre}{\mathcal P}
\newcommand{\Ready}{\mathsf{Ready}}
\newcommand{\Live}{\mathsf{Live}}
\newcommand{\Routes}{\mathsf{Routes}}
\newcommand{\Reach}{\mathsf{Reach}}
\newcommand{\Work}{\mathsf{Work}}
\newcommand{\NoRoute}{\mathsf{NoRoute}}
\newcommand{\Ext}{\mathsf{External}}
\newcommand{\RepCrit}{\mathsf{RepCrit}}
\newcommand{\ResSet}{\mathsf{Res}}
\newcommand{\Par}{\mathsf{Par}}
\newcommand{\Born}{\mathcal Q^{+}}
\newcommand{\ind}[1]{\mathbf 1[#1]}
\newcommand{\res}[2]{\overset{#1}{\Longrightarrow}_{#2}}

\title{Normative Continuity under Self-Revision\\[3pt]
\large How rules can change without erasing authority or obligations}
\author{}
\date{Working synthesis --- 29 August 2026}

\begin{document}
\maketitle

\section*{The question}

Suppose a reasoner can revise the normative rules it represents: its judgments,
warrants, evaluators, adjudicative procedures, and even represented rules about
how those things may be revised. No substantive rule has to remain in force
forever. What has to stay continuous so that revision does not become a way to
create authority from nowhere, erase an inconvenient obligation, change the
rules of a pending case, reset its history, or leave it waiting forever?

The proposal in this note is simple in spirit:
\begin{center}
\emph{the content may change; the history of what was authorized and what became owed may not be silently rewritten.}
\end{center}

There are two parts. \emph{Legitimate Evolution} tracks how rules acquire and
lose authority. The \emph{Answerable Process} tracks what happens to the issues
those rules create. Together they support the headline result, Theorem~\ref{thm:nsa}:
provided no fixed no-route wait can remain unresolved forever, every matter
either closes through recorded resolutions or receives unbounded total
attention.

\paragraph{What is and is not fixed.}
The represented rules may change. The mathematical vocabulary used to read a
history does not: for example, the relation $\Permit_n(c)$ is a fixed part of
the model which asks what the reasoner's \emph{current represented rules} say
about proposed change $c$. In this sense the reasoner may change its represented
rules of revision without changing the metalogical notion of what it means for
those current rules to permit a move. The accepted initial state is also fixed
as a historical starting point. The theory does not claim to eliminate every
fixed meta-level notion.

\paragraph{Kinds of statements.}
A \emph{semantic judgment} says what the reasoner's current represented rules
say at a given history. A \emph{structural requirement} says how the recorded
history must respect those judgments through time. Additional assumptions are
stated only when a theorem needs something neither supplied by the structural
rules nor determined by the represented history.

\section{History and the objects that live in it}\label{sec:objects}

\begin{definition}[history and grounds]
At position $n\in\mathbb N$ the reasoner appends a finite batch of records
\[
 e_n=(r_{n,1},\ldots,r_{n,k_n}),
 \qquad
 H_n=(e_0,\ldots,e_{n-1}).
\]
$H_n$ is the strict prefix available before position $n$. Let
\[
 \Past_n:=\{r_{j,\ell}:j<n\}.
\]
Grounds for a move at position $n$ are simply a finite set of prior records,
\[
 g\fin\Past_n.
\]
Thus a record written in $e_n$ cannot be cited as a prior ground by another
record in the same batch.
\end{definition}

A record may change which normative rules stand in force, open or resolve an
issue, add or withdraw a prerequisite, or mark an issue as a separately
protected matter. Record identity is historical: changing how a later state
interprets the same content does not create a new occurrence.

\subsection{Rules that currently have authority}

At each point in the history, the reasoner treats some represented rules as
currently authoritative. Let $\mathcal N$ be the set of records that can serve
as such rules, and let
\[
 L_n\subseteq\mathcal N
\]
be the rules in force before position $n$. We call a rule in $L_n$
\emph{standing}. The rules in force at the beginning are
\[
 G:=L_0.
\]
$G$ is only the starting point. A rule in $G$ may later be removed, and a rule
not in $G$ may later acquire standing.

Two facts about a standing rule will matter below:
\begin{itemize}[itemsep=1pt,topsep=2pt]
 \item $\Auth(\lambda)$ means that $\lambda$ may be cited as authority for
       changing which rules have standing.
 \item $\lambda\triangleright(\kappa,\tau,x)$ means that $\lambda$ authorizes
       protocol $\kappa$ to handle a new issue of kind $\tau$ about subject $x$.
\end{itemize}
These are features of what the represented rule says. They are not additional
semantic judgments that vary separately with time.

\begin{definition}[current authority to govern an issue]
A protocol is authorized to govern a new issue when some rule already in force
says that it may. We write
\[
 \kappa\Vdash_n(\tau,x)
 \quad\Longleftrightarrow\quad
 \exists\lambda\in L_n:\ \lambda\triangleright(\kappa,\tau,x).
\]
Thus $\kappa\Vdash_n(\tau,x)$ is derived from the rules standing at $H_n$; it
is not a separate source of authority.
\end{definition}

\subsection{Issues}

Let $\mathcal Q$ be the set of issue occurrences. An issue $q$ is the occurrence
of an issue-opening record. It has immutable opening data
\[
 \pi(q)=(\tau_q,x_q,\kappa_q,s_q^0)
 \in \mathcal T\times\mathcal X\times\mathcal K\times\mathcal S_{\kappa_q}.
\]
The four coordinates answer four different questions:
\begin{itemize}[itemsep=1pt,topsep=2pt]
 \item $\tau_q$ --- \emph{what kind of question is this?}
 \item $x_q$ --- \emph{what is the question about?}
 \item $\kappa_q$ --- \emph{which protocol will judge this particular episode?}
 \item $s_q^0$ --- \emph{what local memory does that protocol start with?}
\end{itemize}
The subject $x_q$ may itself be a rule, reason, issue, protocol, event, or
process component. If $q$ is still open at $H_n$, write
$s_n(q)\in\mathcal S_{\kappa_q}$ for its current local state.

Let $\Born_n\subseteq\mathcal Q$ be the issues opened in batch $e_n$; the
superscript $+$ is meant to suggest ``new at $n$.'' The sets $\Born_n$ are
disjoint. Thus
\[
 I_n:=\bigcup_{j<n}\Born_j
\]
is the set of issues incurred before position $n$, and $O_n\subseteq I_n$ is
the set still outstanding at $H_n$.

\subsection{Prerequisites}

Let $\mathcal D$ be the set of prerequisite occurrences. A prerequisite
occurrence is identified with the record that first introduces it, so the same
occurrence cannot later be ``re-added'' after withdrawal. Its immutable data are
\[
 \delta(d)=(q_d,\chi_d,T_d).
\]
Here $q_d$ is the issue that owns the prerequisite, $\chi_d$ is the condition
that would satisfy it, and $T_d\fin\mathcal Q$ is a finite set of issue
occurrences named as initial internal ways of working toward $\chi_d$.
Different prerequisites of one issue are conjunctive; the alternatives in
$T_d$ are disjunctive.

For $q\in O_n$, let $\Pre_n(q)\subseteq\mathcal D$ be the finite set of
prerequisites currently active for $q$.

\section{What the reasoner's current rules decide}\label{sec:semantics}

The following are the \emph{only} primitive prefix-dependent semantic
judgments in this note. Each is evaluated from $H_n$ and may read the current
standing state $L_n$ and whatever represented rules are present in $H_n$.
Everything introduced later is derived from the history plus these judgments.

\begin{semantics}[permission to change standing]
$\Permit_n(c)$ means that the current represented rules permit proposed
standing-change record $c$ at $H_n$.
\end{semantics}

\begin{semantics}[what is currently due]
The current rules determine a set
\[
 \Due_n:=\Due(H_n,L_n)\subseteq\mathcal T\times\mathcal X.
\]
What creates a new obligation is a rising edge,
\[
 \NewDue_n:=
 \begin{cases}
   \Due_0,&n=0,\\
   \Due_n\setminus\Due_{n-1},&n>0.
 \end{cases}
\]
If a pair later drops out of $\Due_n$, that does not close any issue already
opened for it. If it later reappears, that is a new rising edge and can create
a new issue.
\end{semantics}

\begin{semantics}[resolution]
For $q\in O_n$, $g\fin\Past_n$, and finite $S\subseteq\mathcal Q$,
$\Resolve_n(q;g,S)$ means that $q$'s own anchored protocol $\kappa_q$, reading
$H_n$ and $s_n(q)$, accepts disposing of $q$ into successor set $S$ on grounds
$g$. $S=\varnothing$ is a terminal resolution. If batch $e_n$ actually records
such an accepted resolution, write
\[
 q\res{g}{n}S,
 \qquad
 \ResSet_n:=\{q\in O_n:\exists g,S\;q\res{g}{n}S\}.
\]
\end{semantics}

\begin{semantics}[state continuation]
For finite parent set $P\subseteq O_n$, proposed fresh successor $q'$, and
grounds $g$, $\Continue_n(P;q',g)$ means that the current rules accept the
proposed initial state $s_{q'}^0$ as a continuation of the parents' current
states. This relation can allow ordinary state transport or an explicit reset;
what matters structurally is that the transition is represented and accepted,
not silent.
\end{semantics}

\begin{semantics}[prerequisite satisfaction]
$\Met_n(d)$ means that the condition $\chi_d$ is satisfied at prefix $H_n$.
This judgment applies to every prerequisite, not only to information expected
from outside the reasoner.
\end{semantics}

\begin{semantics}[adding and withdrawing prerequisites]
$\AddPre_n(q;g,d)$ means that the current rules accept introducing fresh
prerequisite occurrence $d$ for issue $q$ (so $q_d=q$); $\DropPre_n(q;g,d)$
means that they accept withdrawing active prerequisite $d$ from $q$. Let
$\Pre_n^+(q)$ and $\Pre_n^-(q)$ be the prerequisite occurrences actually added
and withdrawn by accepted records in $e_n$; every $d\in\Pre_n^+(q)$ has
$q_d=q$.
\end{semantics}

\begin{semantics}[creating an additional matter]
$\Designate_n(q;g)$ means that the current rules accept making issue $q$ the
root of an additional protected matter. The issue may already exist or may be
opened in the same batch.
\end{semantics}

For emphasis, the complete list is
\[
 \Permit,\ \Due,\ \Resolve,\ \Continue,\ \Met,\ \AddPre,\ \DropPre,\ \Designate.
\]
Relations such as $\Vdash_n$, readiness, successor ancestry, routes,
reachability, available work, and service are defined from the record and are
not further semantic inputs.

\paragraph{Acceptance convention.}
For the semantic judgments that govern recorded moves, an \emph{accepted} move
means both that its record occurs in $e_n$ and that the corresponding judgment
holds at the strict prefix $H_n$: $\Permit$ for standing changes, $\Resolve$
for resolutions, $\AddPre$/$\DropPre$ for prerequisite changes, and
$\Designate$ for matter designation. $\Continue$ is checked when a successor
opens; $\Due$ and $\Met$ describe the current semantic state rather than
recorded moves.

\section{How authority changes: Legitimate Evolution}\label{sec:standing}

The rules in force are not fixed. At position $n$ the reasoner may add rules to
$L_n$ or remove rules from it. Legitimate Evolution imposes one continuity
condition on this freedom: a change in authority must answer to authority that
was already in force before the change.

Concretely, a proposed change must pass the current rules' own test for whether
it is permitted. And if the change actually adds or removes authority, it must
cite at least one rule that was already standing and was itself allowed to
justify changes of authority.

\begin{requirement}[changing which rules have authority]\label{req:standing-change}
Let $L_n^+$ be the rules that gain standing in batch $e_n$, and let $L_n^-$ be
the rules that lose it. Every record $c$ in $e_n$ that proposes such a change
must satisfy
\[
 \Permit_n(c).
\]
If $c$ actually adds or removes standing, then
\[
 \varnothing\neq g(c)\subseteq\{\lambda\in L_n:\Auth(\lambda)\}.
\]
After the batch, the rules in force are exactly
\[
 L_{n+1}=(L_n\setminus L_n^-)\cup L_n^+.
\]
As in the frozen kernel, the occurrences in $L_n^+$ are fresh: each is first
recorded in $e_n$. A rule that has lost standing does not re-enter as the same
occurrence; readopting its content produces a new occurrence with its own
grounds.
\end{requirement}

Because grounds must come from $\Past_n$, every cited authorizer predates the
change it justifies. In particular, a rule first added in $e_n$ cannot justify
its own admission or another change in the same batch. There is no same-batch
bootstrap of authority.

This local requirement has a global consequence. Pick any rule that is in force
now and ask what earlier standing rules authorized its admission. Ask the same
question about those authorizers, and continue backward. Each step moves to an
earlier position, so this history cannot continue forever.

\begin{theorem}[Grounded Replay]\label{thm:replay}
Assume Requirement~\ref{req:standing-change}. For every $\lambda\in L_n$,
there is a finite authorization tree rooted at $\lambda$ whose leaves lie in
$G$. Each non-leaf node is linked to the standing authorizers cited when that
rule entered $L$, and positions strictly decrease along every branch.
\end{theorem}

\begin{proof}
Induct on the position at which a rule gained standing. A rule not in $G$
entered through a change grounded in earlier standing authorizers, each of
which is in $G$ or entered at a strictly earlier position. Apply the same
argument to each of those grounds. Since positions strictly decrease, every
branch eventually reaches $G$. Freshness of $L_n^+$ makes ``the change through
which the rule entered'' unique.
\end{proof}

Grounded Replay is about where current authority came from, not about keeping
old rules forever. A rule in the ancestry---even a rule in $G$---may later be
legitimately removed. What survives revision is the historical chain by which
present authority was acquired.

This grounding requirement is specific to changes of authority. The structural
theory does not require every resolution, prerequisite change, or matter
designation to cite nonempty grounds. Those moves may have empty grounds unless
their own semantic judgment rules that out. Grounded Replay is therefore a
theorem about the history of authority, not a universal grounding theorem for
every recorded move.

\section{Continuity of issues}\label{sec:issues}

The next requirements say what happens once the current rules make something
due and an issue has been opened for it.

\begin{requirement}[standing at opening]\label{req:anchor}
Every fresh issue uses a protocol that currently has standing for its kind and
subject:
\[
 q\in\Born_n\quad\Longrightarrow\quad
 \kappa_q\Vdash_n(\tau_q,x_q).
\]
A standing change recorded in the same batch cannot authorize the opening:
standing is read from the strict prefix $H_n$.
\end{requirement}

\begin{requirement}[new demands become issues]\label{req:due}
Every newly due pair is represented by a fresh issue:
\[
 (\tau,x)\in\NewDue_n
 \quad\Longrightarrow\quad
 \exists q\in\Born_n:\ (\tau_q,x_q)=(\tau,x).
\]
\end{requirement}

\begin{assumption}[compatibility of what is due with what has standing]\label{assump:compat}
Whenever $(\tau,x)\in\NewDue_n$, at least one protocol currently has standing
for that pair:
\[
 (\tau,x)\in\NewDue_n
 \quad\Longrightarrow\quad
 \exists\kappa\in\mathcal K:\ \kappa\Vdash_n(\tau,x).
\]
This is not a theorem of the structural bookkeeping. It is the minimal
compatibility needed for the previous two requirements to be jointly
satisfiable.
\end{assumption}

\begin{requirement}[resolution continuity]\label{req:resolution-continuity}
\[
 O_{n+1}=(O_n\setminus\ResSet_n)\cup\Born_n.
\]
So an old issue remains outstanding unless it is explicitly resolved, and
every fresh issue is outstanding at the next prefix. In particular,
\[
 q\in O_n\setminus O_{n+1}\quad\Longrightarrow\quad q\in\ResSet_n.
\]
A change in the surrounding normative rules cannot by itself make an issue
disappear.
\end{requirement}

\begin{requirement}[successors are fresh]\label{req:fresh-successors}
If $q\res{g}{n}S$, then $S\subseteq\Born_n$. Define the successor relation
\[
 q\to q'
 \quad\Longleftrightarrow\quad
 \exists n,g,S\;\bigl(q\res{g}{n}S\ \wedge\ q'\in S\bigr).
\]
Because successors open later than their parents, $\to$ is acyclic. If
$q'\in\Born_n$, let
\[
 \Par_n(q'):=\{q\in O_n:\exists g,S\;(q\res{g}{n}S\wedge q'\in S)\}
\]
be the set of parents that resolve into it in that batch.
\end{requirement}

\begin{lemma}[revision is prospective]\label{lem:prospective}
If $q\to q'$, then the resolution of $q$ is judged under $\kappa_q$, while
$q'$ is governed by its own opening anchor $\kappa_{q'}$, which must stand when
$q'$ opens.
\end{lemma}

\begin{proof}
The resolution record for $q$ can be accepted only through
$\Resolve_n(q;g,S)$, whose definition uses the immutable anchor $\kappa_q$. By
Requirement~\ref{req:fresh-successors}, $q'$ is a fresh issue in the same
batch, and Requirement~\ref{req:anchor} checks its own anchor $\kappa_{q'}$
against the standing state already present in $H_n$. Thus revision can govern
the fresh successor without changing the adjudicative terms of the old issue.
\end{proof}

Thus a protocol can gain or lose standing without changing the rules of a case
that is already open. To move a continuing issue to new terms, the process must
resolve the old issue and open a fresh successor under those terms.

\begin{requirement}[state continuity]\label{req:state}
Every fresh successor with at least one parent begins from an explicitly
accepted continuation of the parents' current states:
\[
 q'\in\Born_n,\ \Par_n(q')\neq\varnothing
 \quad\Longrightarrow\quad
 \exists g\fin\Past_n:\ \Continue_n(\Par_n(q');q',g).
\]
The current rules may allow a reset, but the reset must be represented as the
accepted transition to the successor; accumulated local state cannot simply be
replaced without such a transition.
\end{requirement}

\section{Prerequisites and matters}\label{sec:matters}

\subsection{Prerequisites persist until something happens to them}

\begin{requirement}[prerequisite continuity]\label{req:pre-continuity}
If $q$ remains outstanding across position $n$, then
\[
 \Pre_{n+1}(q)
 =\bigl(\Pre_n(q)\setminus\Pre_n^-(q)\bigr)\cup\Pre_n^+(q).
\]
If $q\in\Born_n$ is fresh, then
\[
 \Pre_{n+1}(q)=\Pre_n^+(q).
\]
Thus an active prerequisite of a continuing issue remains active unless an
accepted withdrawal removes it. Resolving the issue ends its remaining
prerequisites with that episode. Because prerequisite identities are fresh
record occurrences, a withdrawn $d$ cannot later be reactivated; a later
prerequisite is a new occurrence $d'$.
\end{requirement}

\begin{requirement}[prerequisites cannot name future issues]\label{req:pre-refs}
If $d$ is first introduced in $e_n$, then
\[
 q_d\in O_n\cup\Born_n,
 \qquad
 T_d\subseteq O_n\cup\Born_n.
\]
A prerequisite may point to issues already open or co-opened in the same
batch, but not to an issue occurrence that does not yet exist.
\end{requirement}

\begin{requirement}[satisfaction is persistent]\label{req:met}
\[
 \Met_n(d)\quad\Longrightarrow\quad\Met_{n+1}(d).
\]
A fixed prerequisite may become satisfied; it cannot later become unsatisfied.
\end{requirement}

\begin{definition}[ready issue]
\[
 \Ready_n(q)
 \quad\Longleftrightarrow\quad
 \forall d\in\Pre_n(q),\ \Met_n(d).
\]
\end{definition}

\begin{requirement}[only ready issues may resolve]\label{req:ready}
\[
 q\in\ResSet_n\quad\Longrightarrow\quad\Ready_n(q).
\]
\end{requirement}

\subsection{A matter is the obligation that survives issue replacement}

An \emph{issue} is one case or episode. A \emph{matter} is the continuing
obligation to deal with a case and whatever successor issues replace it.
This distinction is what prevents procedural churn from resetting service.

Let $\preceq$ be the reflexive-transitive closure of $\to$. Matters are created
in time, not retroactively. Let $M_n$ be the matters that already exist at
prefix $H_n$, with $M_0=\varnothing$. After batch $e_n$, add
\begin{itemize}[itemsep=1pt,topsep=2pt]
 \item every fresh issue $q\in\Born_n$ with $\Par_n(q)=\varnothing$ (a root
       issue), and
 \item every issue $q\in O_n\cup\Born_n$ for which $e_n$ contains an accepted
       $\Designate_n(q;g)$ record.
\end{itemize}
No matter is removed from $M_n$. Write
\[
 M_\infty:=\bigcup_n M_n,
 \qquad
 \beta(m):=\min\{n:m\in M_n\}
\]
for the eventual set of matters and the first prefix at which matter $m$
exists.

\begin{definition}[live descendants of a matter]
For $n\ge\beta(m)$,
\[
 \Live_n(m):=\{q\in O_n:m\preceq q\}.
\]
These are the currently open issues through which the matter continues.
\end{definition}

\begin{lemma}[matter continuity]\label{lem:matter-continuity}
If $n\ge\beta(m)$ and $\Live_n(m)=\varnothing$, then
$\Live_k(m)=\varnothing$ for every $k\ge n$. More generally, every outstanding
descendant of $m$ remains part of matter $m$, through branching and merging.
\end{lemma}

\begin{proof}
An existing descendant belongs to $\Live_n(m)$ by definition. A new descendant
can only be opened as a successor of an outstanding descendant. If no such
outstanding descendant exists, no later one can be created.
\end{proof}

\section{Internal routes and when work is available}\label{sec:work}

A prerequisite can block an issue while pointing to other issues that are
supposed to help satisfy it. The next definitions follow those pointers through
issue succession.

\begin{definition}[current internal routes]
\[
 \Routes_n(d):=\{r\in O_n:\exists t\in T_d\text{ with }t\preceq r\}.
\]
So if a route root $t\in T_d$ is replaced by successors, those live successors
remain routes for the same prerequisite automatically.
\end{definition}

\begin{lemma}[route loss is permanent]\label{lem:route-extinction}
Let $d$ be introduced in $e_j$. If $n>j$ and $\Routes_n(d)=\varnothing$, then
$\Routes_k(d)=\varnothing$ for all $k\ge n$.
\end{lemma}

\begin{proof}
By Requirement~\ref{req:pre-refs}, $T_d\subseteq O_j\cup\Born_j\subseteq I_n$
for $n>j$, so no root in $T_d$ opens at or after $n$. Any route at $k+1\ge n+1$
is therefore either a route at $k$ that remained outstanding, or a fresh issue
in $\Born_k$ that is a successor of an outstanding descendant of some root in
$T_d$; that parent lies in $O_k$ and is itself a route at $k$. Induct on $k$
from $\Routes_n(d)=\varnothing$.
\end{proof}

The restriction $n>j$ is needed: a root $t\in T_d$ co-opened in $e_j$ is not
outstanding at $H_j$, so $\Routes_j(d)$ may be empty while
$\Routes_{j+1}(d)\ni t$. Since $d\in\Pre_n(q_d)$ only for $n>j$, this does not
affect the waiting relation below.

For unmet prerequisites, define a directed waiting relation
\[
 q\rightsquigarrow_n r
 \quad\Longleftrightarrow\quad
 \exists d\in\Pre_n(q):\ \neg\Met_n(d)\ \wedge\ r\in\Routes_n(d).
\]
Read $q\rightsquigarrow_n r$ as: ``$q$ is currently waiting through $r$.''

\begin{definition}[issues relevant to a matter]
$\Reach_n(m)$ is the least set containing $\Live_n(m)$ and closed under
$\rightsquigarrow_n$. In words: start with the live descendants of $m$ and
follow every chain of unmet prerequisites to the internal issues they point to.
At every finite prefix this set is finite.
\end{definition}

\begin{definition}[available work]
\[
 \Work_n(m)
 :=\{q\in\Reach_n(m):\Ready_n(q)\}
 \cup
 \{q\in\Reach_n(m):q\text{ lies on a }\rightsquigarrow_n\text{-cycle}\}.
\]
A ready reachable issue is work that can proceed toward resolution. A waiting
cycle also counts as work: circular internal dependence is itself something
the process can act on, not a reason for permanent idleness.
\end{definition}

Define the opportunity indicator and cumulative opportunity by
\[
 o_n(m):=\ind{m\in M_n\ \wedge\ \Work_n(m)\neq\varnothing},
 \qquad
 \Omega_N(m):=\sum_{n<N}o_n(m).
\]
Thus $\Omega_N(m)$ is simply the number of positions before $N$ at which
matter $m$ had at least one available piece of work.

Let $a_n(m)$ be the service attention assigned to matter $m$ at position $n$.
We require
\[
 0\le a_n(m)\le o_n(m),
 \qquad
 \sum_{m\in M_n}a_n(m)\le 1,
 \qquad
 A_N(m):=\sum_{n<N}a_n(m).
\]
The first inequality says attention is charged only when the matter has
available work; the second is a unit total service budget. Here attention is a
service quantity, not a restriction on which structural records may be written.
For example, the process may record withdrawal of a dead prerequisite while
$o_n(m)=0$; that record can restore future work without itself counting as
attention to currently available work.

\begin{requirement}[non-starvation]\label{req:nonstarvation}
For every $m\in M_\infty$,
\[
 \Omega_N(m)\to\infty
 \quad\Longrightarrow\quad
 A_N(m)\to\infty.
\]
If a matter offers work infinitely often, it cannot receive only finite total
attention.
\end{requirement}

\begin{lemma}[non-starvation is feasible under the unit budget]\label{lem:nonstarve-feasible}
Enumerate matters in order of birth as $m_0,m_1,\ldots$ and give $m_j$ fixed
share $w_j=2^{-j-1}$. Set
\[
 a_n(m_j):=w_j o_n(m_j).
\]
Then $\sum_j a_n(m_j)\le1$ at every position and
$A_N(m_j)=w_j\Omega_N(m_j)$. Hence all matters can satisfy non-starvation
simultaneously. This is only a feasibility witness, not a proposed priority
rule.
\end{lemma}

The final structural rule prevents a process from making a starving matter's
waiting chain move forever.

\begin{requirement}[do not rewire a starving waiting chain]\label{req:no-rewire}
Suppose $q$ is already outstanding, remains outstanding across position $n$,
and receives at least one fresh prerequisite in $e_n$. Then every matter that
already reaches $q$ must currently have some available work:
\[
 q\in O_n\cap O_{n+1},\ \Pre_n^+(q)\neq\varnothing
 \quad\Longrightarrow\quad
 \forall m\in M_n\,
 \bigl(q\in\Reach_n(m)\Rightarrow\Work_n(m)\neq\varnothing\bigr).
\]
Fresh issues may open with prerequisites. The restriction applies only to
adding another prerequisite to an existing issue while some already-existing
matter is stuck through that issue.
\end{requirement}

\begin{remark}[why the rule follows reachability, not ownership]
It is not enough to look only at matters for which $q$ is a direct descendant.
Another matter may be waiting on $q$ through a chain of prerequisites. If that
indirectly waiting matter has no available work, continually replacing $q$'s
prerequisites can keep the matter idle forever while ensuring that no single
prerequisite persists. Requirement~\ref{req:no-rewire} blocks exactly that
failure mode.
\end{remark}

\section{The main structural theorem}\label{sec:main}

\begin{lemma}[every outstanding issue belongs to a matter]\label{lem:root-matter}
For every $q\in O_n$, there is some $m\in M_n$ with $q\in\Live_n(m)$.
\end{lemma}

\begin{proof}
Follow parent links backward from $q$. Parents open earlier than successors, so
the chain is finite and reaches a root issue. That root became a matter when it
opened.
\end{proof}

\begin{lemma}[a reachable resolution creates an opportunity]\label{lem:resolution-opportunity}
If $q\in\Reach_n(m)$ and $q\in\ResSet_n$, then $o_n(m)=1$.
\end{lemma}

\begin{proof}
A resolving issue must be ready by Requirement~\ref{req:ready}. A ready issue
in $\Reach_n(m)$ belongs to $\Work_n(m)$.
\end{proof}

\begin{definition}[a no-route wait]
For $m\in M_n$, let
\[
 \NoRoute_n(m):=
 \left\{d:\begin{array}{l}
 d\in\Pre_n(q)\text{ for some }q\in\Reach_n(m),\\
 \neg\Met_n(d),\quad\Routes_n(d)=\varnothing
 \end{array}\right\}.
\]
So $d\in\NoRoute_n(m)$ means that somewhere along matter $m$'s waiting
structure there is an unmet prerequisite, still active, for which no live
internal issue remains as a route toward satisfaction.
\end{definition}

\begin{theorem}[Persistent-Wait Theorem]\label{thm:persistent-wait}
Assume Requirements \ref{req:resolution-continuity}, \ref{req:fresh-successors},
\ref{req:pre-continuity}--\ref{req:ready} and \ref{req:no-rewire}.
Fix $m\in M_\infty$. If, after its birth,
$m$ remains live forever but has available work only finitely often,
\[
 \exists n_0\ge\beta(m)\ \forall n\ge n_0:\ \Live_n(m)\neq\varnothing,
 \qquad
 \sup_N\Omega_N(m)<\infty,
\]
then eventually one fixed no-route wait remains:
\[
 \exists d,N_0\ \forall n\ge N_0:\ d\in\NoRoute_n(m).
\]
\end{theorem}

\begin{proof}
Since $m\in M_n$ for all $n\ge\beta(m)$ and $\Omega_N(m)$ is bounded, there are
only finitely many positions at which $\Work_n(m)\neq\varnothing$. Choose
$n_1\ge n_0$ after the last of them. Then $\Work_n(m)=\varnothing$ and
$\Live_n(m)\neq\varnothing$ for every $n\ge n_1$.

\emph{1. No reachable issue resolves after $n_1$.}
If $q\in\Reach_n(m)\cap\ResSet_n$ with $n\ge n_1$, Lemma~\ref{lem:resolution-opportunity}
gives $\Work_n(m)\neq\varnothing$, a contradiction.

\emph{2. The reachable waiting structure can only shrink:}
for $n\ge n_1$, $\Live_{n+1}(m)=\Live_n(m)$,
$\Reach_{n+1}(m)\subseteq\Reach_n(m)$, and for every
$q\in\Reach_{n+1}(m)$, $\Pre_{n+1}(q)\subseteq\Pre_n(q)$.

\emph{Live lineage.} Let $q\in\Live_{n+1}(m)$. If $q\in O_n$ then
$q\in\Live_n(m)$. Otherwise $q\in\Born_n$ by
Requirement~\ref{req:resolution-continuity}; $q\neq m$ because $m\in M_n$ was
opened before $n$; so $m\preceq q$ passes through a parent
$p\in\Par_n(q)\subseteq O_n$ with $m\preceq p$. Then $p\in\Live_n(m)$ resolves
at $n$, contradicting Step~1. Conversely nothing in $\Live_n(m)$ resolves, so
$\Live_n(m)\subseteq O_{n+1}$.

\emph{Prerequisites.} Let $q\in\Reach_n(m)$. Then $q\in O_n\cap O_{n+1}$ by
Step~1, and Requirement~\ref{req:no-rewire} with $\Work_n(m)=\varnothing$
gives $\Pre_n^+(q)=\varnothing$, so
$\Pre_{n+1}(q)=\Pre_n(q)\setminus\Pre_n^-(q)\subseteq\Pre_n(q)$ by
Requirement~\ref{req:pre-continuity}.

\emph{Routes.} Let $q\in\Reach_n(m)$, $d\in\Pre_n(q)$ unmet at $n$, and
$r\in\Routes_{n+1}(d)$. Since $d$ was introduced before $n$,
Requirement~\ref{req:pre-refs} puts $T_d\subseteq I_n$, so $r\notin T_d$ if
$r\in\Born_n$. If $r\in O_n$ then $r\in\Routes_n(d)$, hence $q\rightsquigarrow_n r$
and $r\in\Reach_n(m)$ by closure. If $r\in\Born_n$, its parent
$p\in\Par_n(r)\subseteq O_n$ satisfies $t\preceq p$ for some $t\in T_d$, so
$p\in\Routes_n(d)$, $p\in\Reach_n(m)$, and $p$ resolves at $n$, contradicting
Step~1. Thus $\Routes_{n+1}(d)\subseteq\Routes_n(d)$, and since no route in
$\Reach_n(m)$ resolves, in fact $\Routes_{n+1}(d)=\Routes_n(d)$.

\emph{Closure.} Now let $q\in\Reach_{n+1}(m)$ be obtained from
$\Live_{n+1}(m)=\Live_n(m)\subseteq\Reach_n(m)$ by an edge
$q_0\rightsquigarrow_{n+1}q$ with $q_0\in\Reach_n(m)$ (induction along the
closure). The edge is witnessed by some $d\in\Pre_{n+1}(q_0)\subseteq\Pre_n(q_0)$
with $\neg\Met_{n+1}(d)$, hence $\neg\Met_n(d)$ by
Requirement~\ref{req:met}, and $q\in\Routes_{n+1}(d)\subseteq\Routes_n(d)$.
So $q_0\rightsquigarrow_n q$ and $q\in\Reach_n(m)$.

Hence, after $n_1$, all reachable issues lie in the finite set
$\Reach_{n_1}(m)$ and all their active prerequisites in the finite set
$\bigcup_{q\in\Reach_{n_1}(m)}\Pre_{n_1}(q)$; no fresh issue and no fresh
prerequisite occurrence can enter.

\emph{3. The finite structure eventually stabilizes.}
For $n\ge n_1$ the sets $\Reach_n(m)$ and $\Pre_n(q)$ ($q$ reachable) are
non-increasing, and $\Met_n(d)$ is non-decreasing by Requirement~\ref{req:met},
all within one fixed finite set. Route sets of unmet reachable prerequisites
are constant by Step~2. Monotone sequences in a finite set are eventually
constant, so there is $n_2\ge n_1$ after which the reachable issues, active
prerequisites, satisfaction values, and route sets do not change. The waiting
relation $\rightsquigarrow_n$ restricted to $\Reach_n(m)$ is determined by
these, so it is constant too.

\emph{4. A stable finite waiting graph with no available work has a no-route
wait.}
Fix $n\ge n_2$. The matter is still live, so $\Reach_n(m)$ is nonempty. No
reachable issue is ready and no reachable issue lies on a
$\rightsquigarrow_n$-cycle (a self-loop $q\rightsquigarrow_n q$ counts as a
cycle), since either would belong to $\Work_n(m)$. Because $\Reach_n(m)$ is
closed under $\rightsquigarrow_n$, the waiting graph on $\Reach_n(m)$ is a
finite acyclic digraph and has a sink $q^*$. It is not ready, so some active
$d\in\Pre_n(q^*)$ is unmet. If $\Routes_n(d)\neq\varnothing$, $q^*$ would have
an outgoing waiting edge, contrary to being a sink. Thus
$\Routes_n(d)=\varnothing$ and $d\in\NoRoute_n(m)$.

\emph{5. The same wait persists.}
For every $n\ge n_2$ the same $q^*$ is reachable, the same $d$ is active on it
and unmet, and $\Routes_n(d)=\varnothing$, all by stabilization (or, for the
route set, by Lemma~\ref{lem:route-extinction}). Hence $d\in\NoRoute_n(m)$ for
all $n\ge n_2$.
\end{proof}

The theorem deliberately stops here. Structural continuity can show that
permanent idleness must settle on a fixed unmet wait with no internal route. It
does not say why such a wait would eventually clear. For the next result we
need only one liveness assumption; we do not need to classify waits as
internal or external.

\section{From a fixed wait to No Structural Abandonment}\label{sec:nsa}

\begin{assumption}[wait responsiveness]\label{assump:wait}
No fixed prerequisite may remain forever as an active, reachable, unmet wait
with no internal route. Formally, for every matter $m$, prerequisite
occurrence $d$, and position $N_0$,
\[
 \bigl(\forall n\ge N_0:\ d\in\NoRoute_n(m)\bigr)
 \quad\Longrightarrow\quad
 \exists k\ge N_0:\ \Met_k(d).
\]
This deliberately abstracts over why the wait clears. A prerequisite that is
currently no-route may instead be withdrawn or cease to be reachable before it
becomes permanent. Coverage of outside inputs and a separate discipline for
removing obsolete internal prerequisites are possible sufficient conditions for
this one assumption; the structural theory itself uses only the combined
liveness statement above.
\end{assumption}

\begin{theorem}[Persistent Opportunity]\label{thm:persistent-opportunity}
Under the hypotheses of Theorem~\ref{thm:persistent-wait} and wait
responsiveness, if matter $m$ remains live forever after its birth, then
\[
 \Omega_N(m)\to\infty.
\]
\end{theorem}

\begin{proof}
If cumulative opportunity were bounded, Theorem~\ref{thm:persistent-wait}
would produce $d$ and $N_0$ with $d\in\NoRoute_n(m)$ for every $n\ge N_0$.
Wait responsiveness gives some $k\ge N_0$ with $\Met_k(d)$, contradicting
$d\in\NoRoute_k(m)$, which by definition requires $\neg\Met_k(d)$.
\end{proof}

\begin{theorem}[No Structural Abandonment]\label{thm:nsa}
Add Requirement~\ref{req:nonstarvation} (non-starvation). Then for every matter
$m\in M_\infty$, at least one of the following holds (they need not be
exclusive):
\[
 \bigl(\exists N\ge\beta(m)\ \forall n\ge N:\ \Live_n(m)=\varnothing\bigr)
 \qquad\text{or}\qquad
 A_N(m)\to\infty.
\]
In words: the matter eventually closes through explicit resolutions, or it
receives unbounded total attention.
\end{theorem}

\begin{proof}
After a matter is born, emptiness of $\Live_n(m)$ is permanent by
Lemma~\ref{lem:matter-continuity}. Thus either the matter eventually closes or
it remains live forever. In the second case Theorem~\ref{thm:persistent-opportunity}
gives $\Omega_N(m)\to\infty$, and non-starvation gives $A_N(m)\to\infty$.
Resolution continuity makes the first case explicit rather than silent.
\end{proof}

\begin{center}
\emph{If no fixed no-route wait can persist forever: resolve it, or keep servicing it.}
\end{center}

\section{One continuity principle, repeated at several levels}\label{sec:picture}

The individual rules are easier to remember as instances of one idea:
revision may change normative content, but it may not silently break the
historical relation that carries authority or obligation forward.

\begin{center}
\begin{tabular}{@{}p{0.27\textwidth}p{0.62\textwidth}@{}}
\toprule
\textbf{Part of the theory} & \textbf{What it prevents} \\
\midrule
Grounded Replay & A currently standing rule appearing without an authorization history. \\
New demands become issues & A demand generated by the reasoner's own rules remaining merely implicit. \\
Resolution continuity & An outstanding issue disappearing when the surrounding rules change. \\
Prospective revision & A new rule retroactively becoming the judge of an already-open case. \\
State continuity & A successor silently forgetting the procedural state accumulated by its parents. \\
Matter continuity & Renaming, splitting, or replacing issues resetting the obligation to deal with the underlying matter. \\
Persistent-Wait Theorem & Internal restructuring producing permanent idleness without eventually exposing one fixed unmet wait. \\
Non-starvation & Infinitely many opportunities for work receiving only finite total service. \\
\bottomrule
\end{tabular}
\end{center}

Two short slogans capture the two halves:
\[
 \boxed{\text{authority must have a history}}
 \qquad\qquad
 \boxed{\text{obligations must have a future}}.
\]

\section{How this connects to criticism and substantive progress}\label{sec:updown}

The structural theory starts only once something is represented. It does not
show that important criticism enters the record, and it does not show that
attention leads to a good answer. Those are separate assumptions or theorems.

A discovery/coverage theory might supply:
\begin{quote}
If a relevant criticism remains exposed in the world, eventually a record
representing that criticism appears.
\end{quote}
The reasoner's own current normative rules might then satisfy a responsiveness
condition such as
\[
 \RepCrit_n(c)
 \quad\Longrightarrow\quad
 (\mathrm{Registration},c)\in\NewDue_{n+1}.
\]
The one-step shift is just strict-prefix timing: a criticism written in $e_n$
first becomes readable in $H_{n+1}$. Coverage of outside inputs can also help
justify wait responsiveness; removing obsolete internal prerequisites supplies
another possible source of that assumption. The present theory does not need to
distinguish the two cases.

At the other end, an application-specific progress theorem might say
\[
 A_N(m)\to\infty
 \quad+\quad
 \text{persistent relevant evidence}
 \quad\Longrightarrow\quad
 \text{a useful accepted disposition}.
\]
Nothing in the present structural theory says the critic must win. The claim is
only that a protected matter cannot be erased, silently re-judged, silently
reset, hidden forever behind a moving chain of prerequisites, or starved when
work keeps becoming available.

The full intended composition is therefore
\[
\begin{aligned}
&\text{persistent relevant criticism}
  \longrightarrow \text{represented criticism}
  \longrightarrow \text{newly due issue}
  \longrightarrow \text{matter},\\[4pt]
&\text{matter}
  \xrightarrow{\ \text{wait responsiveness}\ } \text{explicit closure or infinitely much opportunity},\\[4pt]
&\text{infinitely much opportunity}
  \longrightarrow \text{unbounded attention}
  \longrightarrow \text{useful disposition}.
\end{aligned}
\]
Each arrow needs its own assumption or theorem.

\section{Scope and proof status}\label{sec:scope}

\paragraph{What this theory does not establish.}
It does not establish normative correctness, good judgment, reflective
openness, a final evaluator, rates of progress, or resistance to manipulating
which future evidence becomes available. A self-sealing rule can have a
perfectly legitimate authorization history. Structural continuity says how
such a rule must be changed and how the obligations it creates persist; it does
not say the rule is substantively good.

\paragraph{What remains fixed.}
The accepted initial standing state $G$ and the eight semantic relations of
Section~\ref{sec:semantics} are fixed parts of the mathematical model. Their
values may depend on the current represented rules, and those represented rules
may themselves change. Internalizing still more of this meta-level structure is
a separate problem.

\paragraph{Current provenance and proof status.}
The standing-change discipline and Grounded Replay restate the frozen
Legitimate Evolution work. Resolution continuity and the basic issue/state
continuity rules consolidate the subsequent Answerable Process design.
The route-following construction, the no-rewiring rule, and
Theorems~\ref{thm:persistent-wait}--\ref{thm:nsa} are the newer part of the
synthesis. The proofs here are paper proofs and should remain marked as proof
obligations until independently checked and entered into the project
concordance. In particular, the no-rewiring rule is motivated by the explicit
cross-matter prerequisite-rotation counterexample that breaks the weaker
ownership-only gate. The notation $\Work_n(m)$ is the same mathematical object
previously called the actionable frontier $\mathsf{Front}_n(m)$; this is an
expository rename, not a change in the definition.

\paragraph{Which requirements each theorem uses.}
Grounded Replay uses Requirement~\ref{req:standing-change} only. The
Persistent-Wait Theorem uses Requirements \ref{req:resolution-continuity},
\ref{req:fresh-successors}, \ref{req:pre-continuity}--\ref{req:ready} and
\ref{req:no-rewire}, together with finiteness of batches; Persistent
Opportunity adds wait responsiveness, and No Structural Abandonment adds
Requirement~\ref{req:nonstarvation}. Requirements \ref{req:anchor},
\ref{req:due} and \ref{req:state}, the compatibility assumption, and the
judgments $\Permit$, $\Due$, $\Continue$ and $\Designate$ are part of the
theory of normative continuity but are not load-bearing for any theorem after
Grounded Replay; the standing layer and the issue layer meet only at
Requirement~\ref{req:anchor}. Lemma~\ref{lem:root-matter} is likewise not used
by the theorems.

\paragraph{Adjacent work is not superseded.}
This synthesis does not replace the separate Proper Exercise results about
jurisdictional self-ratification and the absence of a generic no-escalation
theorem; the checker/trace-agreement layer that asks whether a concrete run
actually follows the represented semantics; or the evidence/uptake/answerability
decomposition developed for Legitimate Improvement. Those theories answer
different questions and remain separate inputs or downstream consumers.

\paragraph{Audit status.}
\textsc{agent-consolidated} (29 August 2026): the definitions and theorems
above were independently reconstructed and adversarially checked by an agent,
with executable fixtures in \texttt{normative\_continuity\_fixtures.py}
(the rotating-prerequisite countermodel against the ownership-only gate, its
rejection by Requirement~\ref{req:no-rewire}, co-opened route roots, route
extinction, waiting cycles, and branching/merging/designation). This label
does not mean \textsc{frozen}, \textsc{canonical}, \textsc{proved}, or
\textsc{lean-verified}. The next stages are provenance concordance and
independent or formal proof checking.

\paragraph{Next mathematical debt.}
The priced and clocked charge/liability machinery should appear only through a
separate theorem showing that a concrete bounded reasoner satisfies these
structural requirements, and stating whatever stronger finite-time or
liability guarantees that construction adds. It should not be built into the
structural definitions themselves.

\end{document}
